%\chapter{Erdschluss in Mittelspannungsnetzen}

\chapter{Erdschlussberechnung und Erdschlussverhalten}
\label{chap:09-erdschluss}

\section{Einordnung}

Ein Erdschluss liegt vor, wenn ein aktiver Leiter über eine Fehlerstelle
mit Erde oder einem geerdeten leitfähigen Teil verbunden wird. Der
Erdschlussstrom hängt nicht nur von der Netzspannung und der Fehlerimpedanz
ab, sondern wesentlich von der Sternpunktbehandlung, den kapazitiven
Strömen, den Nullsystemimpedanzen und den Erdungswegen.

Die Beurteilung eines Erdschlusses muss deshalb zwischen unterschiedlichen
Netzformen unterscheiden. Insbesondere sind isolierte, kompensierte und
niederohmig geerdete Netze nicht mit demselben vereinfachten Modell
zu behandeln.

\section{Grundlegende Netzformen}

\subsection{Isoliertes Netz}

In einem isolierten Netz ist der Netzsternpunkt nicht absichtlich niederohmig
mit Erde verbunden. Ein einpoliger Erdschluss wird daher wesentlich durch
die Kapazitäten der gesunden Phasen gegen Erde und die übrigen
Ableitpfade bestimmt.

Der Erdschlussstrom kann bei ausgedehnten Kabelnetzen trotz fehlender
direkter Sternpunkterdung beträchtlich sein, da die Kabelkapazitäten
einen kapazitiven Rückstrompfad bilden.

\subsection{Kompensiertes Netz}

In einem kompensierten Netz wird der Sternpunkt über eine
Petersen-Spule mit Erde verbunden. Die induktive Stromkomponente der Spule
wird so eingestellt, dass sie den kapazitiven Erdschlussstrom des Netzes
weitgehend kompensiert.

Die Kompensation ist nicht automatisch vollständig. Netzänderungen,
Schaltzustände, Verluste und Verstimmungen beeinflussen den verbleibenden
Erdschlussstrom. Die tatsächliche Auslegung und Abstimmung müssen zum
Netz und zum vorgesehenen Schutzkonzept passen.

\subsection{Niederohmig geerdetes Netz}

In niederohmig geerdeten Netzen ist der Sternpunkt über eine definierte
Impedanz, häufig einen Widerstand, geerdet. Der Erdschlussstrom wird
dadurch gezielt begrenzt und zugleich ein ausreichend deutliches
Schutzsignal ermöglicht.

Der Widerstand beeinflusst die Höhe des Fehlerstroms, die
Sternpunktverschiebung und die Erdungsspannungen. Die Auslegung muss
Fehlerstrom, zulässige Fehlerdauer, Schutzansprechen und thermische
Beanspruchung gemeinsam berücksichtigen.

\section{Einpoliger Erdschluss und symmetrische Komponenten}

Ein einpoliger Erdschluss an Phase \(a\) kann unter üblichen Voraussetzungen
mit den Sequenznetzen beschrieben werden. Für einen Fehler mit Impedanz
\(Z_f\) gelten am Fehlerpunkt
\[
  I_b=0,\qquad I_c=0.
\]

Daraus folgt für die Sequenzströme
\[
  I_0=I_1=I_2.
\]
Der Fehlerstrom beträgt
\[
  I_f=I_a=I_0+I_1+I_2=3I_0.
\]

Die Sequenznetze werden bei diesem Fehler in Reihe geschaltet. Mit der
Mitsystem-Ersatzspannung \(U_{\mathrm{Th}}\) ergibt sich im vereinfachten
Modell
\[
  I_f=
  \frac{3U_{\mathrm{Th}}}
  {Z_1+Z_2+Z_0+3Z_f}.
\]

Diese Gleichung ist nur dann unmittelbar anwendbar, wenn die verwendeten
Sequenzimpedanzen und die Ersatzspannung die konkrete Netzform abbilden.
Insbesondere darf die Nullsystemimpedanz nicht ohne Prüfung durch die
Mitsystemimpedanz ersetzt werden.

\section{Bedeutung der Nullsystemimpedanz}

Die Nullsystemimpedanz hängt unter anderem ab von:
\begin{itemize}
  \item der Sternpunktbehandlung;
  \item Transformator-Schaltgruppen und Wicklungsverschaltungen;
  \item Erdungsimpedanzen und Erdungsleitern;
  \item Kabelschirmen und metallischen Bewehrungen;
  \item Rückstromwegen über Erde und geerdete leitfähige Strukturen;
  \item dem Aufbau der Netzkomponenten.
\end{itemize}

Eine Erdschlussberechnung, die nur die Mitsystemimpedanz berücksichtigt,
kann deshalb erheblich vom tatsächlichen Verhalten abweichen.

\section{Sternpunktverschiebung und gesunde Phasen}

Bei einem einpoligen Erdschluss in einem isolierten oder kompensierten Netz
kann sich das Potential des Netzsternpunkts gegenüber Erde verschieben.
Dadurch verändern sich die Spannungen der Phasen gegen Erde.

In einem idealisierten isolierten Netz kann die Spannung der beiden
gesunden Phasen gegen Erde bei einem metallischen einpoligen Erdschluss
auf ungefähr die verkettete Spannung ansteigen. Die tatsächlichen
Spannungen hängen von Netzkapazitäten, Dämpfung, Fehlerimpedanz und
Transienten ab.

Diese Spannungsverschiebung ist für die Isolationsbeanspruchung der
gesunden Phasen wichtig. Sie ist ein Grund dafür, dass die Netzform und
das Erdschlusskonzept bei der Auswahl der Isolationskoordination
berücksichtigt werden müssen.

\section{Kapazitiver Erdschlussstrom}

Für eine einzelne Kapazität \(C\) gegen Erde gilt bei sinusförmiger Spannung
\[
  I_C=\omega C U.
\]

Für ein symmetrisches Drehstromnetz ist bei der Berechnung des gesamten
Erdschlussstroms sorgfältig zu unterscheiden, welche Phasenkapazitäten
und welche Phasenspannungen in den jeweiligen Strompfad eingehen.

Für ein idealisiertes isoliertes Netz mit gleicher Phasenkapazität \(C\)
gegen Erde kann der kapazitive Erdschlussstrom eines metallischen
einpoligen Erdschlusses näherungsweise als
\[
  I_E\approx3\omega C U_{\mathrm{ph}}
\]
angegeben werden, wobei \(U_{\mathrm{ph}}\) die normale Strangspannung
gegen Erde und \(C\) die Kapazität einer Phase gegen Erde bezeichnet.

Die Gleichung setzt voraus, dass die drei Phasenkapazitäten gleich sind
und dass die betrachteten Rückstrompfade durch das kapazitive Modell
hinreichend beschrieben werden. Bei realen Kabelnetzen können zusätzliche
Kapazitäten, Schirme, Erdungsverbindungen und Verluste relevant sein.

\subsection{Beispiel}

Ein idealisiertes 20-kV-Netz habe eine gesamte wirksame Phasenkapazität
gegen Erde von
\[
  C=1\,\mathrm{\mu F}
\]
je Phase. Bei \(f=50\,\mathrm{Hz}\) ist
\[
  \omega=2\pi f\approx314{,}16\,\mathrm{rad/s}.
\]

Die Strangspannung beträgt
\[
  U_{\mathrm{ph}}=\frac{20\,\mathrm{kV}}{\sqrt{3}}
  \approx11{,}55\,\mathrm{kV}.
\]

Damit ergibt sich im idealisierten Modell
\[
  I_E\approx3\cdot314{,}16\cdot1\,\mathrm{\mu F}
  \cdot11{,}55\,\mathrm{kV}
  \approx10{,}9\,\mathrm A.
\]

Dieser Wert beschreibt nur den kapazitiven Anteil unter den angegebenen
Annahmen. Er ist nicht ohne Weiteres der gesamte Fehlerstrom eines
realen Netzes.

\section{Kompensation durch Petersen-Spule}

Die Petersen-Spule stellt eine induktive Verbindung zwischen Sternpunkt
und Erde her. Ihre Stromkomponente ist näherungsweise
\[
  I_L=\frac{U_{\mathrm{ph}}}{\omega L}
\]
für eine ideale Induktivität \(L\), an der die Spannung \(U_{\mathrm{ph}}\)
anliegt.

Die Abstimmung zielt darauf, den induktiven Spulenstrom an den kapazitiven
Erdschlussstrom des Netzes anzupassen. In einem vereinfachten Modell
kann der Resonanzzustand durch die Gleichheit der Beträge der
kapazitiven und induktiven Anteile beschrieben werden.

Reale Netze enthalten Verluste und nichtideale Elemente. Deshalb bleibt
ein Wirkstromanteil, und eine vollständige Kompensation ist nicht
automatisch gewährleistet. Zusätzlich muss die Regelung auf
Netzumschaltungen und Änderungen der wirksamen Kapazität reagieren können.

\section{Erdschlussimpedanz und Fehlerort}

Der Erdschluss ist nicht notwendigerweise metallisch. Die Fehlerstelle
kann eine erhebliche Impedanz besitzen, beispielsweise durch Lichtbogen,
verschmutzte Isolatoren, feuchte Oberflächen oder unvollständigen Kontakt
mit Erde.

Eine erhöhte Fehlerimpedanz reduziert häufig den Fehlerstrom. Dadurch kann
die sichere Erkennung durch einen Überstromschutz erschwert werden.
Erdschlussschutzkonzepte nutzen deshalb je nach Netzform unter anderem
Residualstrom, Nullspannung, Richtungsinformation oder weitere
netzabhängige Messgrößen.

Der Fehlerstrom allein genügt nicht in jedem Netz, um Fehlerort und
Fehlerrichtung sicher zu bestimmen.

\section{Schutztechnische Erkennung}

Die Erkennung eines Erdschlusses richtet sich nach dem Netzkonzept.
Mögliche Messgrößen sind:
\begin{itemize}
  \item Summenstrom beziehungsweise Residualstrom;
  \item Nullsystemstrom \(I_0\);
  \item Nullsystemspannung \(U_0\);
  \item Richtungsinformation aus Strom und Spannung;
  \item transiente oder wattmetrische Komponenten, abhängig vom Schutzprinzip.
\end{itemize}

Die konkrete Auswahl hängt von der Sternpunktbehandlung, der Fehlerimpedanz,
dem Umfang des Netzes und den verfügbaren Messwandlern ab. Die
Einstellwerte sind anhand der tatsächlichen Netzparameter und der
Schutzkoordination zu bestimmen.

\section{Erdungsspannung und Berührungsspannung}

Bei Erdschlüssen kann Strom in Erdungsanlagen und den umgebenden Boden
fließen. Dadurch entstehen Potentialgradienten. Zu unterscheiden sind:
\begin{itemize}
  \item Erdpotentialanhebung einer Erdungsanlage;
  \item Berührungsspannung zwischen gleichzeitig berührbaren leitfähigen Teilen;
  \item Schrittspannung zwischen zwei Punkten auf der Erdoberfläche;
  \item über leitfähige Verbindungen eingetragene Potentiale.
\end{itemize}

Die Höhe der Erdpotentialanhebung allein reicht nicht aus, um die
Personengefährdung zu beurteilen. Entscheidend sind unter anderem die
Potentialverteilung, Strompfade durch den Körper, Fehlerdauer,
Oberflächenbeschaffenheit und die konkrete Erdungsgeometrie.

Die Bewertung und Auslegung müssen nach den einschlägigen Normen und
den konkreten Anlagenbedingungen erfolgen. Ein vereinfachtes
Ersatzschaltbild ersetzt keine Erdungsberechnung.

\section{Vergleich der Netzformen}

\begin{center}
\begin{tabular}{p{0.22\linewidth}p{0.34\linewidth}p{0.34\linewidth}}
\hline
\textbf{Netzform} & \textbf{Prägendes Verhalten} & \textbf{Wichtige Prüfgröße}\\
\hline
Isoliert &
Erdschlussstrom wird wesentlich durch Netzkapazitäten und Ableitpfade bestimmt &
Kapazitiver Erdschlussstrom und Sternpunktverschiebung\\
\hline
Kompensiert &
Induktiver Spulenstrom kompensiert den kapazitiven Anteil weitgehend &
Verstimmung, Reststrom und Schutzempfindlichkeit\\
\hline
Niederohmig geerdet &
Sternpunkterdung begrenzt und bestimmt den Erdschlussstrom gezielt &
Fehlerstrom, Fehlerdauer und thermische Beanspruchung\\
\hline
\end{tabular}
\end{center}

Die Tabelle beschreibt allgemeine Tendenzen. Sie ersetzt keine konkrete
Berechnung der jeweiligen Netzform.

\section{Rechenablauf für einen einpoligen Erdschluss}

Ein systematischer Ablauf umfasst:
\begin{enumerate}
  \item Netzform und Sternpunktbehandlung identifizieren.
  \item Fehlerstelle, Fehlerphase und Fehlerimpedanz definieren.
  \item Mit-, Gegen- und Nullsystemimpedanzen ermitteln.
  \item Transformator-Schaltgruppen, Kabelschirme und Erdungswege berücksichtigen.
  \item Die Sequenznetze gemäß den Fehlerbedingungen verschalten.
  \item Den Fehlerstrom und die Nullsystemgrößen berechnen.
  \item Spannungsverschiebung und Schutzgrößen bewerten.
  \item Bei Bedarf Erdpotentialanhebung, Berührungs- und Schrittspannungen
        mit einem geeigneten Erdungsmodell untersuchen.
  \item Ergebnisse gegen die tatsächlichen Anlagenparameter und die
        geltenden Normen prüfen.
\end{enumerate}

\section{Grenzen der vereinfachten Rechnung}

Die analytische Erdschlussrechnung ist nur so gut wie das verwendete
Nullsystemmodell. Insbesondere bei ausgedehnten Kabelnetzen, komplexen
Erdungsanlagen, Lichtbogenfehlern oder leistungselektronischen Einspeisern
kann ein einfaches konzentriertes Ersatzmodell unzureichend sein.

Bei transienten Erdschlüssen spielen zusätzlich Schaltvorgänge, Wanderwellen,
nichtlineare Lichtbogenkennlinien und die frequenzabhängigen Eigenschaften
der Netzkomponenten eine Rolle. Solche Vorgänge erfordern gegebenenfalls
eine zeitbereichsorientierte Simulation.

\section{Übungsaufgaben}

\begin{enumerate}
  \item Erklären Sie, weshalb die Nullsystemimpedanz bei einem Erdschluss
        nicht ohne Weiteres aus der Mitsystemimpedanz abgeleitet werden kann.
  \item Berechnen Sie für ein idealisiertes isoliertes 10-kV-Netz mit
        \(C=0{,}8\,\mathrm{\mu F}\) je Phase gegen Erde den kapazitiven
        Erdschlussstrom bei \(50\,\mathrm{Hz}\).
  \item Beschreiben Sie den Unterschied zwischen isoliertem,
        kompensiertem und niederohmig geerdetem Netz.
  \item Warum kann eine hohe Fehlerimpedanz die Erdschlusserkennung erschweren?
  \item Erläutern Sie, weshalb die Erdpotentialanhebung allein keine
        vollständige Bewertung der Berührungsspannung ermöglicht.
  \item Welche Parameter müssen bekannt sein, bevor eine quantitative
        Erdschlussberechnung für ein reales Mittelspannungsnetz durchgeführt
        werden kann?
\end{enumerate}